\documentclass[12pt]{article}
\usepackage[utf8]{inputenc}
\usepackage[T1]{fontenc}
\usepackage[margin=1in]{geometry}
\usepackage{setspace}
\usepackage{xcolor}
\usepackage[colorlinks=true, urlcolor=blue!60!black, linkcolor=blue!60!black]{hyperref}
\setlength{\parindent}{2em}
\setlength{\parskip}{1.25em}
\renewcommand{\baselinestretch}{1.0}
\begin{document}
\begin{center}
{\LARGE \textit{Washington v.\ Davis}, 426 U.S.\ 229 (1976)\par}
\vspace{0.8em}
{\large Case Brief --- Four-Page Edition\par}
\end{center}
\vspace{1em}

\section{Memory Jogger}

A facially neutral policy with a racially disproportionate impact does not violate equal protection---under the Fifth or the Fourteenth Amendment---unless the challenger proves a racially discriminatory purpose (\href{https://supreme.justia.com/cases/federal/us/426/229/}{Justia}).

\section{Detailed Case Facts}

In 1970 two Black police officers, including Alfred Davis, sued District of Columbia officials over promotion policies in the Metropolitan Police Department. Two rejected Black applicants, George Harley and John Sellers, intervened to challenge recruiting under the Fifth Amendment's Due Process Clause (which carries an equal protection component under \textit{Bolling v.\ Sharpe}), 42 U.S.C.\ \S~1981, and the D.C.\ Code. Recruits had to score at least 40 of 80 on ``Test 21,'' a Civil Service Commission examination of verbal ability, vocabulary, reading, and comprehension. The intervenors did not allege intentional discrimination; they argued the test was unrelated to job performance and that four times as many Black as white applicants failed it. But since 1969, 44\% of new recruits had been Black, the Department had affirmatively recruited Black candidates, and a validation study showed Test 21 scores correlated with recruit-training results.

\section{Procedural History}

The district court granted summary judgment for the officials: the recruitment record and the test's link to training performance negated any inference of discrimination. Over a dissent, the D.C.\ Circuit reversed, importing the Title VII disparate-impact rule of \textit{Griggs v.\ Duke Power Co.} into the Constitution and holding the test invalid unless shown to predict job performance. The Supreme Court reversed on June 7, 1976.

\section{Judicial Votes}

\textbf{7--2}, reversed.
\begin{itemize}
\item \textbf{Majority:} Justice White, joined by Chief Justice Burger and Justices Blackmun, Powell, Rehnquist, and Stevens; Justice Stewart joined Parts I and II.
\item \textbf{Concurring:} Justice Stevens.
\item \textbf{Dissenting:} Justice Brennan, joined by Justice Marshall (on the statutory holding).
\end{itemize}

\section{Holding}

A facially neutral law is not unconstitutional solely because it has a racially disproportionate impact; a challenger must show discriminatory purpose to trigger strict scrutiny, though impact is relevant circumstantial evidence. Absent such purpose, rational basis applies. Under the applicable statutes, a test may be validated by its relationship to performance in a job-relevant training program. Justice White's warning about tax, welfare, and licensing statutes and footnote 12's disapproval of lower-court cases were influential dicta.

\section{Analysis of Opinions}

\textbf{Majority (White).} Noticing plain error, the Court held that the D.C.\ Circuit wrongly collapsed the constitutional standard into Title VII. Jury-exclusion, apportionment, school-desegregation (\textit{Keyes}), and welfare (\textit{Jefferson v.\ Hackney}) cases all required purpose. Purpose may be inferred from ``the totality of the relevant facts,'' including severe impact as in \textit{Yick Wo}. The government may legitimately seek to upgrade verbal skills; Black and white applicants who failed stood on the same footing; and the recruiting record negated intent. A disparate-impact rule would threaten a vast array of neutral laws, so any expansion is for legislatures. Part III upheld training-program validation under the statutes.

\textbf{Stevens (concurring).} Purpose and impact are not sharply distinct: the ``most probative evidence of intent'' is often objective, since actors are presumed to intend natural consequences. But ``[a] law conscripting clerics should not be invalidated because an atheist voted for it.'' Literacy is relevant to policing, and the test was used nationwide.

\textbf{Brennan, joined by Marshall (dissenting).} The Court should not have reached the statutory question, which was not properly presented, and footnote 12 prejudged \textit{Arlington Heights}. On the merits, \textit{Griggs} and \textit{Albemarle} job-relatedness standards were unmet: correlating one verbal test with another risks rewarding verbal facility, not ``job-specific ability.''

\section{Examples: Future Applications}

\begin{enumerate}
\item \textbf{Same-Side: Licensing Fee.} A doubled barber-licensing fee that disproportionately burdens minority entrepreneurs, adopted without discriminatory intent, gets rational basis review.
\item \textbf{Same-Side: Veterans' Preference.} A lifetime veterans' preference that overwhelmingly benefits men survives; foreseeable impact is not purpose (\textit{Feeney}, 1979).
\item \textbf{Opposite-Side: Rezoning.} A suburb's abrupt rezoning after an integrated project is proposed, with procedural departures and coded remarks, shows purpose under \textit{Arlington Heights}.
\item \textbf{Opposite-Side: Selective Permits.} Granting 95\% of white applicants' permits while denying 98\% of Asian American applicants supports an inference of intent, as in \textit{Yick Wo}.
\item \textbf{Fence-Sitter: Risk Algorithm.} A bail algorithm relying on race-correlated proxies, adopted despite warnings, leans toward validity under \textit{Davis}, but Stevens's objective-evidence approach supports a purpose argument.
\end{enumerate}

\section{Critique}

\textbf{Antisubordination critics.} Charles Lawrence argued the intent rule ignores unconscious racism; Alan Freeman faulted its ``perpetrator perspective''; Reva Siegel described an ``equality divided'' that defers to neutral policies while strictly scrutinizing race-conscious remedies. The rule also invites officials to cloak motives in neutral language.

\textbf{Defenders.} Paul Brest and originalists such as Earl Maltz and John Harrison argue that equal protection forbids purposeful denial of equal legal status, not unequal outcomes, and that a disparate-impact rule would shift policymaking to judges while Congress remains free to legislate impact remedies.

\section{Key Quotations}

\begin{quote}
``Disproportionate impact is not irrelevant, but it is not the sole touchstone of an invidious racial discrimination forbidden by the Constitution. Standing alone, it does not trigger the rule, \textit{McLaughlin v.\ Florida}, 379 U.\ S.\ 184 (1964), that racial classifications are to be subjected to the strictest scrutiny, and are justifiable only by the weightiest of considerations.''
\end{quote}
\noindent --- Justice White, for the Court, 426 U.S.\ at 242.

\begin{quote}
``Frequently the most probative evidence of intent will be objective evidence of what actually happened, rather than evidence describing the subjective state of mind of the actor. For, normally, the actor is presumed to have intended the natural consequences of his deeds.''
\end{quote}
\noindent --- Justice Stevens, concurring, 426 U.S.\ at 253.

\begin{quote}
``The central purpose of the Equal Protection Clause of the Fourteenth Amendment is the prevention of official conduct discriminating on the basis of race. It is also true that the Due Process Clause of the Fifth Amendment contains an equal protection component prohibiting the United States from invidiously discriminating between individuals or groups. \textit{Bolling v.\ Sharpe}, 347 U.\ S.\ 497 (1954). But our cases have not embraced the proposition that a law or other official act, without regard to whether it reflects a racially discriminatory purpose, is unconstitutional solely because it has a racially disproportionate impact.''
\end{quote}
\noindent --- Justice White, for the Court, 426 U.S.\ at 239.

\end{document}
